% =====================================================================
% Note on: Bond Implied CDS Spread and CDS-Bond Basis
% Richard Zhou, 2008  —  ssrn:1265548
% =====================================================================
\documentclass[11pt,a4paper]{article}

\newcommand{\papertag}{zhou\_2008\_cds}

% =====================================================================
% Shared preamble for paper notes
% Include with:  % =====================================================================
% Shared preamble for paper notes
% Include with:  % =====================================================================
% Shared preamble for paper notes
% Include with:  \input{../_common/preamble}
% =====================================================================

% --- Core packages ---
\usepackage[utf8]{inputenc}
\usepackage[T1]{fontenc}
\usepackage{lmodern}
\usepackage[margin=2.5cm]{geometry}
\usepackage{amsmath,amssymb,amsthm,mathtools}
\usepackage{bm}
\usepackage{booktabs}
\usepackage{array}
\usepackage{enumitem}
\usepackage{hyperref}
\usepackage{xcolor}
\usepackage{listings}
\usepackage{fancyhdr}
\usepackage{titlesec}

% --- Hyperref styling ---
\hypersetup{
  colorlinks=true,
  linkcolor=blue!50!black,
  urlcolor=blue!50!black,
  citecolor=blue!50!black
}

% --- Theorem-like environments ---
\theoremstyle{definition}
\newtheorem{defn}{Definition}
\newtheorem{remark}{Remark}
\newtheorem{example}{Example}

\theoremstyle{plain}
\newtheorem{prop}{Proposition}
\newtheorem{lemma}{Lemma}

% --- Coloured boxes for the structured sections ---
% Validation block, commentary, TODOs use coloured frames so the eye
% can find them when flipping through the note.
\usepackage[most]{tcolorbox}

\newtcolorbox{commentary}{
  colback=yellow!5,
  colframe=yellow!50!black,
  title=Commentary,
  fonttitle=\bfseries,
  breakable
}

\newtcolorbox{validation}{
  colback=green!3,
  colframe=green!50!black,
  title=Validation,
  fonttitle=\bfseries,
  breakable
}

\newtcolorbox{todobox}{
  colback=red!3,
  colframe=red!50!black,
  title=TODO / Open questions,
  fonttitle=\bfseries,
  breakable
}

% --- Code listing style for Python snippets in validation blocks ---
\lstdefinestyle{py}{
  language=Python,
  basicstyle=\ttfamily\footnotesize,
  keywordstyle=\color{blue!60!black},
  commentstyle=\color{gray},
  stringstyle=\color{red!60!black},
  showstringspaces=false,
  breaklines=true,
  frame=single,
  framesep=4pt,
  rulecolor=\color{gray!30}
}

% --- Page header: shows paper tag so a printed note is identifiable ---
\pagestyle{fancy}
\fancyhf{}
\fancyhead[L]{\small\textit{\papertag}}
\fancyhead[R]{\small\thepage}
\renewcommand{\headrulewidth}{0.4pt}

% --- Section spacing: tighter than article default ---
\titlespacing*{\section}{0pt}{1.5ex plus 0.5ex minus 0.2ex}{1ex plus 0.2ex}
\titlespacing*{\subsection}{0pt}{1.2ex plus 0.3ex minus 0.2ex}{0.6ex plus 0.2ex}

% =====================================================================

% --- Core packages ---
\usepackage[utf8]{inputenc}
\usepackage[T1]{fontenc}
\usepackage{lmodern}
\usepackage[margin=2.5cm]{geometry}
\usepackage{amsmath,amssymb,amsthm,mathtools}
\usepackage{bm}
\usepackage{booktabs}
\usepackage{array}
\usepackage{enumitem}
\usepackage{hyperref}
\usepackage{xcolor}
\usepackage{listings}
\usepackage{fancyhdr}
\usepackage{titlesec}

% --- Hyperref styling ---
\hypersetup{
  colorlinks=true,
  linkcolor=blue!50!black,
  urlcolor=blue!50!black,
  citecolor=blue!50!black
}

% --- Theorem-like environments ---
\theoremstyle{definition}
\newtheorem{defn}{Definition}
\newtheorem{remark}{Remark}
\newtheorem{example}{Example}

\theoremstyle{plain}
\newtheorem{prop}{Proposition}
\newtheorem{lemma}{Lemma}

% --- Coloured boxes for the structured sections ---
% Validation block, commentary, TODOs use coloured frames so the eye
% can find them when flipping through the note.
\usepackage[most]{tcolorbox}

\newtcolorbox{commentary}{
  colback=yellow!5,
  colframe=yellow!50!black,
  title=Commentary,
  fonttitle=\bfseries,
  breakable
}

\newtcolorbox{validation}{
  colback=green!3,
  colframe=green!50!black,
  title=Validation,
  fonttitle=\bfseries,
  breakable
}

\newtcolorbox{todobox}{
  colback=red!3,
  colframe=red!50!black,
  title=TODO / Open questions,
  fonttitle=\bfseries,
  breakable
}

% --- Code listing style for Python snippets in validation blocks ---
\lstdefinestyle{py}{
  language=Python,
  basicstyle=\ttfamily\footnotesize,
  keywordstyle=\color{blue!60!black},
  commentstyle=\color{gray},
  stringstyle=\color{red!60!black},
  showstringspaces=false,
  breaklines=true,
  frame=single,
  framesep=4pt,
  rulecolor=\color{gray!30}
}

% --- Page header: shows paper tag so a printed note is identifiable ---
\pagestyle{fancy}
\fancyhf{}
\fancyhead[L]{\small\textit{\papertag}}
\fancyhead[R]{\small\thepage}
\renewcommand{\headrulewidth}{0.4pt}

% --- Section spacing: tighter than article default ---
\titlespacing*{\section}{0pt}{1.5ex plus 0.5ex minus 0.2ex}{1ex plus 0.2ex}
\titlespacing*{\subsection}{0pt}{1.2ex plus 0.3ex minus 0.2ex}{0.6ex plus 0.2ex}

% =====================================================================

% --- Core packages ---
\usepackage[utf8]{inputenc}
\usepackage[T1]{fontenc}
\usepackage{lmodern}
\usepackage[margin=2.5cm]{geometry}
\usepackage{amsmath,amssymb,amsthm,mathtools}
\usepackage{bm}
\usepackage{booktabs}
\usepackage{array}
\usepackage{enumitem}
\usepackage{hyperref}
\usepackage{xcolor}
\usepackage{listings}
\usepackage{fancyhdr}
\usepackage{titlesec}

% --- Hyperref styling ---
\hypersetup{
  colorlinks=true,
  linkcolor=blue!50!black,
  urlcolor=blue!50!black,
  citecolor=blue!50!black
}

% --- Theorem-like environments ---
\theoremstyle{definition}
\newtheorem{defn}{Definition}
\newtheorem{remark}{Remark}
\newtheorem{example}{Example}

\theoremstyle{plain}
\newtheorem{prop}{Proposition}
\newtheorem{lemma}{Lemma}

% --- Coloured boxes for the structured sections ---
% Validation block, commentary, TODOs use coloured frames so the eye
% can find them when flipping through the note.
\usepackage[most]{tcolorbox}

\newtcolorbox{commentary}{
  colback=yellow!5,
  colframe=yellow!50!black,
  title=Commentary,
  fonttitle=\bfseries,
  breakable
}

\newtcolorbox{validation}{
  colback=green!3,
  colframe=green!50!black,
  title=Validation,
  fonttitle=\bfseries,
  breakable
}

\newtcolorbox{todobox}{
  colback=red!3,
  colframe=red!50!black,
  title=TODO / Open questions,
  fonttitle=\bfseries,
  breakable
}

% --- Code listing style for Python snippets in validation blocks ---
\lstdefinestyle{py}{
  language=Python,
  basicstyle=\ttfamily\footnotesize,
  keywordstyle=\color{blue!60!black},
  commentstyle=\color{gray},
  stringstyle=\color{red!60!black},
  showstringspaces=false,
  breaklines=true,
  frame=single,
  framesep=4pt,
  rulecolor=\color{gray!30}
}

% --- Page header: shows paper tag so a printed note is identifiable ---
\pagestyle{fancy}
\fancyhf{}
\fancyhead[L]{\small\textit{\papertag}}
\fancyhead[R]{\small\thepage}
\renewcommand{\headrulewidth}{0.4pt}

% --- Section spacing: tighter than article default ---
\titlespacing*{\section}{0pt}{1.5ex plus 0.5ex minus 0.2ex}{1ex plus 0.2ex}
\titlespacing*{\subsection}{0pt}{1.2ex plus 0.3ex minus 0.2ex}{0.6ex plus 0.2ex}

% =====================================================================
% House notation: shared symbols used across all paper notes.
% Each note imports this and may shadow individual macros if a paper
% uses a clashing convention — but the *default* meaning is fixed here.
%
% Include with:  % =====================================================================
% House notation: shared symbols used across all paper notes.
% Each note imports this and may shadow individual macros if a paper
% uses a clashing convention — but the *default* meaning is fixed here.
%
% Include with:  % =====================================================================
% House notation: shared symbols used across all paper notes.
% Each note imports this and may shadow individual macros if a paper
% uses a clashing convention — but the *default* meaning is fixed here.
%
% Include with:  \input{../_common/notation}
% =====================================================================

% --- Measures and expectations ---
\newcommand{\Q}{\mathbb{Q}}              % risk-neutral measure
\newcommand{\Qt}{\mathbb{Q}^{T}}         % T-forward measure
\newcommand{\E}{\mathbb{E}}              % expectation
\newcommand{\Et}[1]{\E^{#1}}             % expectation under measure #1
\newcommand{\Ftime}[1]{\mathcal{F}_{#1}} % filtration at #1

% --- Discount and forward objects ---
\newcommand{\Disc}[2]{D_{#1,#2}}         % discount factor D_{t,T}
\newcommand{\Bond}[2]{P_{#1}(#2)}        % bond price P_t(y)
\newcommand{\Ann}[1]{A_{#1}}             % annuity / level
\newcommand{\Numer}[1]{N_{#1}}           % numeraire
\newcommand{\Fwd}[2]{F_{#1,#2}}          % forward price F_{t,T}
\newcommand{\Fwdpx}[2]{\mathrm{Fwd}_{#1}(#2)} % verbose forward-price F_t(T)

% --- Rates ---
\newcommand{\IRR}{\mathrm{IRR}}          % internal rate of return
\newcommand{\Yld}{y}                     % bond yield variable
\newcommand{\Strike}{K}                  % strike (price-space)
\newcommand{\Lock}{L}                    % locked yield / rate
\newcommand{\Repo}{r^{\mathrm{repo}}}    % repo rate
\newcommand{\Fund}{r^{\mathrm{fun}}}     % unsecured funding rate
\newcommand{\OIS}{r^{\mathrm{ois}}}      % OIS rate
\newcommand{\Coll}{r^{\mathrm{c}}}       % collateral rate
\newcommand{\Rcmt}{R^{\mathrm{cmt}}}     % constant-maturity-treasury rate
\newcommand{\Rswp}{R^{\mathrm{swp}}}     % swap rate (used for risky variant)
\newcommand{\Rec}{\mathrm{Rec}}          % recovery rate
\newcommand{\NPV}{\mathrm{NPV}}          % present-value functional
\newcommand{\PVone}{\mathrm{PV01}}       % risky annuity
\newcommand{\Loss}{\mathrm{Loss}}        % default-leg integral
\newcommand{\Sasw}{S^{\mathrm{ASW}}}     % asset-swap spread
\newcommand{\Scds}{S^{\mathrm{CDS}}}     % CDS spread
\newcommand{\Basis}{\mathrm{Basis}}      % CDS-bond basis

% --- Sensitivities ---
\newcommand{\Diff}[2]{D_{#1}^{#2}}       % Euler's notation: D_x^k[f]
\newcommand{\Riskf}{\mathrm{RiskFactor}} % bond risk factor (= -DV01*1e4)
\newcommand{\DVone}{\mathrm{DV01}}

% --- Indicator / misc ---
\newcommand{\Ind}{\mathbf{1}}
\newcommand{\given}{\,|\,}
\newcommand{\dd}{\,\mathrm{d}}

% --- Greeks-style ---
\newcommand{\Gam}{\Gamma}
\newcommand{\Vega}{\mathcal{V}}

% =====================================================================

% --- Measures and expectations ---
\newcommand{\Q}{\mathbb{Q}}              % risk-neutral measure
\newcommand{\Qt}{\mathbb{Q}^{T}}         % T-forward measure
\newcommand{\E}{\mathbb{E}}              % expectation
\newcommand{\Et}[1]{\E^{#1}}             % expectation under measure #1
\newcommand{\Ftime}[1]{\mathcal{F}_{#1}} % filtration at #1

% --- Discount and forward objects ---
\newcommand{\Disc}[2]{D_{#1,#2}}         % discount factor D_{t,T}
\newcommand{\Bond}[2]{P_{#1}(#2)}        % bond price P_t(y)
\newcommand{\Ann}[1]{A_{#1}}             % annuity / level
\newcommand{\Numer}[1]{N_{#1}}           % numeraire
\newcommand{\Fwd}[2]{F_{#1,#2}}          % forward price F_{t,T}
\newcommand{\Fwdpx}[2]{\mathrm{Fwd}_{#1}(#2)} % verbose forward-price F_t(T)

% --- Rates ---
\newcommand{\IRR}{\mathrm{IRR}}          % internal rate of return
\newcommand{\Yld}{y}                     % bond yield variable
\newcommand{\Strike}{K}                  % strike (price-space)
\newcommand{\Lock}{L}                    % locked yield / rate
\newcommand{\Repo}{r^{\mathrm{repo}}}    % repo rate
\newcommand{\Fund}{r^{\mathrm{fun}}}     % unsecured funding rate
\newcommand{\OIS}{r^{\mathrm{ois}}}      % OIS rate
\newcommand{\Coll}{r^{\mathrm{c}}}       % collateral rate
\newcommand{\Rcmt}{R^{\mathrm{cmt}}}     % constant-maturity-treasury rate
\newcommand{\Rswp}{R^{\mathrm{swp}}}     % swap rate (used for risky variant)
\newcommand{\Rec}{\mathrm{Rec}}          % recovery rate
\newcommand{\NPV}{\mathrm{NPV}}          % present-value functional
\newcommand{\PVone}{\mathrm{PV01}}       % risky annuity
\newcommand{\Loss}{\mathrm{Loss}}        % default-leg integral
\newcommand{\Sasw}{S^{\mathrm{ASW}}}     % asset-swap spread
\newcommand{\Scds}{S^{\mathrm{CDS}}}     % CDS spread
\newcommand{\Basis}{\mathrm{Basis}}      % CDS-bond basis

% --- Sensitivities ---
\newcommand{\Diff}[2]{D_{#1}^{#2}}       % Euler's notation: D_x^k[f]
\newcommand{\Riskf}{\mathrm{RiskFactor}} % bond risk factor (= -DV01*1e4)
\newcommand{\DVone}{\mathrm{DV01}}

% --- Indicator / misc ---
\newcommand{\Ind}{\mathbf{1}}
\newcommand{\given}{\,|\,}
\newcommand{\dd}{\,\mathrm{d}}

% --- Greeks-style ---
\newcommand{\Gam}{\Gamma}
\newcommand{\Vega}{\mathcal{V}}

% =====================================================================

% --- Measures and expectations ---
\newcommand{\Q}{\mathbb{Q}}              % risk-neutral measure
\newcommand{\Qt}{\mathbb{Q}^{T}}         % T-forward measure
\newcommand{\E}{\mathbb{E}}              % expectation
\newcommand{\Et}[1]{\E^{#1}}             % expectation under measure #1
\newcommand{\Ftime}[1]{\mathcal{F}_{#1}} % filtration at #1

% --- Discount and forward objects ---
\newcommand{\Disc}[2]{D_{#1,#2}}         % discount factor D_{t,T}
\newcommand{\Bond}[2]{P_{#1}(#2)}        % bond price P_t(y)
\newcommand{\Ann}[1]{A_{#1}}             % annuity / level
\newcommand{\Numer}[1]{N_{#1}}           % numeraire
\newcommand{\Fwd}[2]{F_{#1,#2}}          % forward price F_{t,T}
\newcommand{\Fwdpx}[2]{\mathrm{Fwd}_{#1}(#2)} % verbose forward-price F_t(T)

% --- Rates ---
\newcommand{\IRR}{\mathrm{IRR}}          % internal rate of return
\newcommand{\Yld}{y}                     % bond yield variable
\newcommand{\Strike}{K}                  % strike (price-space)
\newcommand{\Lock}{L}                    % locked yield / rate
\newcommand{\Repo}{r^{\mathrm{repo}}}    % repo rate
\newcommand{\Fund}{r^{\mathrm{fun}}}     % unsecured funding rate
\newcommand{\OIS}{r^{\mathrm{ois}}}      % OIS rate
\newcommand{\Coll}{r^{\mathrm{c}}}       % collateral rate
\newcommand{\Rcmt}{R^{\mathrm{cmt}}}     % constant-maturity-treasury rate
\newcommand{\Rswp}{R^{\mathrm{swp}}}     % swap rate (used for risky variant)
\newcommand{\Rec}{\mathrm{Rec}}          % recovery rate
\newcommand{\NPV}{\mathrm{NPV}}          % present-value functional
\newcommand{\PVone}{\mathrm{PV01}}       % risky annuity
\newcommand{\Loss}{\mathrm{Loss}}        % default-leg integral
\newcommand{\Sasw}{S^{\mathrm{ASW}}}     % asset-swap spread
\newcommand{\Scds}{S^{\mathrm{CDS}}}     % CDS spread
\newcommand{\Basis}{\mathrm{Basis}}      % CDS-bond basis

% --- Sensitivities ---
\newcommand{\Diff}[2]{D_{#1}^{#2}}       % Euler's notation: D_x^k[f]
\newcommand{\Riskf}{\mathrm{RiskFactor}} % bond risk factor (= -DV01*1e4)
\newcommand{\DVone}{\mathrm{DV01}}

% --- Indicator / misc ---
\newcommand{\Ind}{\mathbf{1}}
\newcommand{\given}{\,|\,}
\newcommand{\dd}{\,\mathrm{d}}

% --- Greeks-style ---
\newcommand{\Gam}{\Gamma}
\newcommand{\Vega}{\mathcal{V}}


\title{Note on: \emph{Bond Implied CDS Spread and CDS-Bond Basis}\\
       \large Richard Zhou, 2008}
\author{Reading notes}
\date{\today}

\begin{document}
\maketitle

% =====================================================================
\section*{Metadata}
\begin{tabular}{@{}p{3.5cm}p{11cm}@{}}
\toprule
Title           & Bond Implied CDS Spread and CDS-Bond Basis \\
Author          & Richard Zhou \\
Email           & \texttt{rzhou50@gmail.com} \\
Date            & 15 August 2008 \\
Source          & \href{https://ssrn.com/abstract=1265548}{ssrn:1265548} \\
Local file      & \texttt{papers/zhou\_2008\_cds/ssrn1265548.pdf} \\
Tags            & \texttt{cds}, \texttt{basis}, \texttt{asset-swap}, \texttt{negative-basis-trade},
                  \texttt{survival}, \texttt{recovery}, \texttt{libor-discounting} \\
Date read       & \today \\
\bottomrule
\end{tabular}

% =====================================================================
\section*{One-paragraph summary}
A short, explicit derivation of the CDS spread implied by a bond's
market price, obtained from the no-arbitrage condition on a negative
basis trade (buy bond, fund at Libor, buy CDS protection on
$1 - D / (1 - R)$ notional --- the amount that makes the combined
position default neutral). The main formula
\eqref{eq:zhou-cds-explicit} expresses the bond-implied CDS as the
sum of three weighted contributions: bond coupon, an
average-Libor-weighted term, and a bond-discount term, all scaled by
risky / risk-free annuities. A consistency theorem proves this CDS
spread also satisfies the standard CDS pricing identity
\eqref{eq:zhou-cds-standard} \emph{if and only if} the survival
probabilities and recovery are consistent with the bond price. The
analogous decomposition for the par-asset-swap spread
\eqref{eq:zhou-asw} and the resulting three-term decomposition of the
CDS-bond basis \eqref{eq:zhou-basis} make the basis structure
transparent: an interest-rate-curve-slope term, a default-treatment
asymmetry term (always negative), and a bond-discount term
(positive for discount bonds, negative for premium bonds). Numerical
tables on a 10Y 7\%-coupon bond confirm the qualitative predictions.

% =====================================================================
\section{Setup and notation}

\begin{tabular}{@{}p{3cm}p{3.5cm}p{8.5cm}@{}}
\toprule
Paper             & House                       & Meaning \\
\midrule
$T_0, T_k$        & $T_0, T_k$                  & Today; bond coupon / ASW payment dates \\
$C$               & ---                         & Bond coupon rate (fixed, semiannual in examples) \\
$M$               & ---                         & Bond risk-free price \\
$1 - D$           & ---                         & Bond market (dirty) price; $D$ is the discount from par \\
$R$               & $\Rec$                      & Recovery rate (40\% in examples) \\
$W$               & $W = (1 - D/(1-R))/(1-D)$   & CDS-notional-to-bond-price ratio (default-neutral) \\
$L_{k-1}$         & $L_{k-1}$                   & Forward Libor for period $(T_{k-1}, T_k)$ at $T_0$ \\
$DF(T_k)$         & $\Disc{T_0}{T_k}$           & Risk-free discount factor \\
$P(\tau > T_k)$   & $Q(T_0, T_k)$               & Survival probability to $T_k$ \\
$\lambda$         & $\lambda$                   & Hazard rate (constant in single-bond calibration) \\
$A$               & $A$                         & Risk-free annuity \\
$\PVone$          & ---                         & Risky annuity (survival-weighted) \\
$\overline{\PVone}$ & ---                       & Weighted-average risky annuity (paper's $\widehat{PV01}$) \\
$\overline L$     & $\overline L$               & DF-weighted average forward Libor (eq.~\ref{eq:zhou-Lbar}) \\
$\overline L^{\,\text{Risky}}$ & ---            & Risky-DF-weighted average forward Libor \\
$\Loss$           & ---                         & Risky default-leg integral $\sum (DF_{k-1} - DF_k)(P_{k-1} + P_k)/2$ \\
$\Sasw$           & ---                         & Par asset swap spread \\
$\Scds$           & ---                         & Bond-implied CDS spread \\
$\Basis$          & ---                         & $\Scds - \Sasw$, CDS-bond basis \\
\bottomrule
\end{tabular}

% =====================================================================
\section{Key equations}

\paragraph{Risk-free annuity and DF-weighted Libor:}
\begin{equation}\label{eq:zhou-A-L}
  A \;=\; \sum_{k=1}^{N} \Delta T_{k-1}\, DF(T_k), \qquad
  \overline L \;=\; \frac{1}{A} \sum_{k=1}^{N} L_{k-1}\, \Delta T_{k-1}\, DF(T_k).
\end{equation}

\paragraph{Par asset swap spread (two equivalent forms):}
\begin{equation}\label{eq:zhou-asw}
  \Sasw \;=\; \frac{M - (1 - D)}{A} \;=\; C - \overline L \;-\; \frac{D}{A}.
\end{equation}
\textit{Says:} the first form is the textbook ``risk-free price minus
market price, amortised by the annuity''; always non-negative when
$M \ge 1 - D$. The second form decomposes the spread into
\emph{coupon $-$ avg Libor $-$ amortised discount}, which the paper
exploits for the basis decomposition. Note: $-D/A$ is positive when
$D > 0$ (discount bond).

\paragraph{Default-neutral CDS notional:}
\begin{equation}\label{eq:zhou-W}
  W \;=\; \frac{1 - D / (1 - R)}{1 - D}.
\end{equation}
\textit{Says:} buying $W \cdot (1 - D)$ of CDS protection makes the
investor's loss on default equal exactly $1 - D$, recovering the
purchase price. For a discount bond ($D > 0$), $W < 1$ (less CDS than
notional); for a premium bond, $W > 1$.

\paragraph{Risky annuity and Loss leg:}
\begin{equation}\label{eq:zhou-pv-loss}
  \PVone \;=\; \sum_{k=1}^{N} \Delta T_{k-1}\, DF(T_k)\, P(\tau > T_k), \qquad
  \Loss \;=\; \sum_{k=1}^{N} \bigl(DF(T_{k-1}) - DF(T_k)\bigr)\, \frac{P(\tau > T_{k-1}) + P(\tau > T_k)}{2}.
\end{equation}

\paragraph{Standard CDS pricing (consistency target):}
\begin{equation}\label{eq:zhou-cds-standard}
  \Scds \;=\; \frac{(1 - R)\, \Loss}{\PVone}.
\end{equation}
\textit{Says:} the textbook CDS-fair-spread formula.

\paragraph{Survival-recovery consistency with bond price:}
\begin{equation}\label{eq:zhou-bond-decomp}
  1 - D \;=\; C \cdot \PVone \;+\; R \cdot \Loss \;+\; P(\tau > T_N) \cdot DF(T_N).
\end{equation}
\textit{Says:} a defaultable bond decomposes into the risky-discounted
coupon stream, the recovery on the default leg, and the survival-
weighted final notional. If $\lambda$ (or a term structure of
$\lambda$) is calibrated so that this identity holds at the observed
bond price, the model is self-consistent.

\paragraph{Bond-implied CDS spread (explicit form):}
\begin{equation}\label{eq:zhou-cds-explicit}
  \Scds \;=\; \frac{C \cdot \PVone}{W \cdot \PVone^{\text{Risky}}}
          \;-\; \frac{\overline L^{\,\text{Risky}} \cdot \PVone}{W \cdot \PVone^{\text{Risky}}}
          \;+\; \frac{D}{W \cdot \PVone^{\text{Risky}}}\!\left(1 - \frac{R}{1 - D} \cdot \frac{\Loss}{\PVone^{\text{Risky}}}\right)\!\!,
\end{equation}
which under \eqref{eq:zhou-bond-decomp} simplifies to (paper's eq.~7)
\begin{equation}\label{eq:zhou-cds-simplified}
  \Scds \;=\; \frac{C \cdot \PVone - \overline L^{\,\text{Risky}} \cdot \PVone^{\text{Risky}}}{\PVone^{\text{Risky}}}
            \;+\; \frac{D}{\PVone^{\text{Risky}}}.
\end{equation}
\textit{Says:} bond-implied CDS = (coupon-vs-Libor differential
weighted by risky annuities) + (bond-discount term per unit risky
annuity). The form makes each input's contribution legible.

\paragraph{CDS-bond basis decomposition} (paper's eq.~8):
\begin{equation}\label{eq:zhou-basis}
  \Basis \;=\; \Scds - \Sasw
          \;=\; \underbrace{\bigl( \overline L - \overline L^{\,\text{Risky}} \bigr)}_{\text{(i) curve-slope}}
          \;-\; \underbrace{C \!\left(1 - \frac{\PVone}{\PVone^{\text{Risky}}}\right)}_{\text{(ii) accrual asymmetry, $\ge 0$ so subtracted}}
          \;+\; \underbrace{D \!\left(\frac{1}{\PVone^{\text{Risky}}} - \frac{1}{A}\right)}_{\text{(iii) off-par}}.
\end{equation}
\textit{Says:} three additive contributions to the basis:
\begin{itemize}[noitemsep]
  \item (i) sign of curve slope: positive when forward Libor is upward
        sloping, negative for inverted, zero for flat.
  \item (ii) bond/CDS accrual asymmetry: always negative
        contribution to the basis since
        $\PVone < \PVone^{\text{Risky}}$? Actually $\PVone^{\text{Risky}} < \PVone$
        when survival is below 1; the paper writes the term so the
        net contribution is negative. (Coupon accrued but unpaid on
        bond default; CDS premium accrued and \emph{paid} on default.)
  \item (iii) off-par effect: positive for discount bonds ($D > 0$),
        negative for premium bonds, vanishing at par.
\end{itemize}

% =====================================================================
\section{Derivations \& commentary}

\begin{commentary}
The cash-flow accounting that produces \eqref{eq:zhou-cds-explicit}
(Appendix A in the paper) is worth internalising: a negative-basis
trade has zero initial value, three streams of running cashflows
(net coupon, accrued premiums on default, terminal redemption
on survival), and the no-arbitrage condition $\mathbb{E}^Q[\cdot] = 0$
applied with the Lando-formula factorisation
$\mathbb{E}^Q[\Ind_{\tau > T} \cdot Z(T)] = DF(T)\, P(\tau > T)$
(under credit-rate independence) gives exactly
\eqref{eq:zhou-cds-explicit}. The trapezoidal approximation of the
default-leg integral produces the $(P_{k-1} + P_k)/2$ weights in
$\Loss$.
\end{commentary}

\begin{commentary}
\textbf{Consistency theorem (Appendix B).} The proposition shows
\eqref{eq:zhou-cds-explicit} $\Rightarrow$ \eqref{eq:zhou-cds-standard}
under \eqref{eq:zhou-bond-decomp}. The key algebraic identity is
$\Loss = (1 - R)\, [\text{tail-recovery term}] + L \cdot \PVone - L^{\text{Risky}} \cdot \PVone^{\text{Risky}}$,
which lets the explicit form be rewritten in terms of $(1 - R)\, \Loss / \PVone$.
The reverse implication does \emph{not} hold: a CDS spread quoted in
the market with an assumed flat recovery $R$ may not satisfy
\eqref{eq:zhou-bond-decomp} on the corresponding bond, which is the
empirical CDS-bond basis.
\end{commentary}

\begin{commentary}
The paper's three application scenarios are clean and worth lifting:
\textbf{(2.3.1)} single bond: assume $R$, solve \eqref{eq:zhou-bond-decomp}
for constant $\lambda$, plug into \eqref{eq:zhou-cds-simplified}.
\textbf{(2.3.2)} bond term structure: bootstrap a piecewise-constant
$\lambda$ across maturities. \textbf{(2.3.3)} both CDS and bond term
structures: solve \eqref{eq:zhou-cds-simplified} for $\lambda$ using
market CDS spreads, then use \eqref{eq:zhou-bond-decomp} to solve
for an implied recovery-rate term structure --- comparable to
recovery-swap quotes.
\end{commentary}

\begin{commentary}
The interpretation of $W$ in \eqref{eq:zhou-W} is operationally
important. A natural alternative is to buy CDS on full $(1-D)$
notional, but this leaves a residual exposure $DR$ on default. The
paper notes (Remark 1 after eq.~7) that the simpler $1-D$ notional
gives the \emph{same} formula \eqref{eq:zhou-cds-simplified} for the
fair spread; the smaller default payoff is offset by larger carry on
the loan. So the choice of CDS notional is presentational, not
mathematical.
\end{commentary}

\begin{commentary}
\textbf{Libor as ``risk-free''.} The paper takes Libor as risk-free
because at the time CDS markets discounted at Libor (this is 2008,
pre-OIS regime). Modern practice would use OIS for both bond
discounting and CDS pricing, and reinstate a separate
funding-vs-discount distinction. The structural formulas
\eqref{eq:zhou-cds-explicit}--\eqref{eq:zhou-basis} are unchanged
under OIS discounting; only the input curve changes.
\end{commentary}

% =====================================================================
\section{Validation}

\begin{validation}
This is a \emph{specification} of mathematical objects the paper
defines and the numerical values they should take. It says nothing
about how those objects are implemented or invoked.

\textbf{Quantities the paper defines:}
\begin{itemize}[noitemsep]
  \item Risk-free annuity $A$ and DF-weighted Libor $\overline L$ \eqref{eq:zhou-A-L}.
  \item Par asset swap spread $\Sasw$ via \eqref{eq:zhou-asw} (two
        equivalent forms).
  \item Default-neutral CDS notional ratio $W$ \eqref{eq:zhou-W}.
  \item Risky annuity $\PVone^{\text{Risky}}$ and default-leg
        integral $\Loss$ \eqref{eq:zhou-pv-loss}.
  \item Bond-pricing consistency identity \eqref{eq:zhou-bond-decomp}.
  \item Bond-implied CDS spread \eqref{eq:zhou-cds-explicit} and its
        simplified form \eqref{eq:zhou-cds-simplified}.
  \item Standard CDS spread \eqref{eq:zhou-cds-standard}.
  \item CDS-bond basis decomposition \eqref{eq:zhou-basis}.
\end{itemize}

\textbf{Canonical numerical case} (Tables 1 \& 2):
\begin{itemize}[noitemsep]
  \item Bond: 10Y, 7\% semiannual coupon, 30/360 day count.
  \item $R = 40\%$.
  \item Libor: swap zero curve for 16-July-2008 from Bloomberg
        (3M $= 2.8\%$ to 10Y $= 4.8\%$, linearly interpolated to
        payment dates). Table 2 uses flat $4.7\%$ as a comparison.
  \item Calibration mode: single-bond (\S2.3.1), constant $\lambda$
        from \eqref{eq:zhou-bond-decomp}.
  \item $D \in \{-10\%, -5\%, 0\%, 5\%, 10\%, 15\%, 20\%\}$ (bond
        market price ranges from $110$ to $80$).
\end{itemize}

\textbf{Expected numerical values} (Table 1, market curve; all in \%):

\smallskip
\begin{tabular}{@{}lrrrrrrr@{}}
\toprule
$D$ (\%)                            & $-10$  & $-5$   & $0$    & $5$    & $10$   & $15$   & $20$  \\
\midrule
$W$                                 & $1.06$ & $1.03$ & $1.00$ & $0.96$ & $0.93$ & $0.88$ & $0.83$ \\
$\Scds$                             & $0.96$ & $1.60$ & $2.31$ & $3.09$ & $3.97$ & $4.97$ & $6.13$ \\
$\Sasw$                             & $1.05$ & $1.67$ & $2.29$ & $2.92$ & $3.54$ & $4.16$ & $4.79$ \\
$\Basis$                            & $-0.09$& $-0.07$& $0.01$ & $0.17$ & $0.43$ & $0.81$ & $1.34$ \\
$\lambda$                           & $1.58$ & $2.64$ & $3.80$ & $5.09$ & $6.65$ & $8.19$ & $10.11$ \\
$\overline L - \overline L^{\text{Risky}}$ & $0.03$ & $0.05$ & $0.08$ & $0.11$ & $0.14$ & $0.17$ & $0.21$ \\
$C\,(1 - \PVone / \PVone^{\text{Risky}})$ & $0.03$ & $0.05$ & $0.07$ & $0.09$ & $0.11$ & $0.14$ & $0.18$ \\
$D / \PVone^{\text{Risky}} - D / A$ & $-0.09$& $-0.07$& $0.00$ & $0.15$ & $0.40$ & $0.78$ & $1.31$ \\
\bottomrule
\end{tabular}
\smallskip

\textit{Table 2 (flat $4.7\%$ Libor curve, same bond):}

\smallskip
\begin{tabular}{@{}lrrrrrrr@{}}
\toprule
$D$ (\%)                            & $-10$  & $-5$   & $0$    & $5$    & $10$   & $15$   & $20$  \\
\midrule
$\Scds$                             & $0.92$ & $1.55$ & $2.24$ & $3.00$ & $3.86$ & $4.83$ & $5.95$ \\
$\Sasw$                             & $1.04$ & $1.67$ & $2.30$ & $2.93$ & $3.56$ & $4.20$ & $4.83$ \\
$\Basis$                            & $-0.12$& $-0.12$& $-0.06$& $0.07$ & $0.29$ & $0.63$ & $1.12$ \\
$\lambda$                           & $1.51$ & $2.55$ & $3.68$ & $4.94$ & $6.35$ & $7.96$ & $9.81$ \\
\bottomrule
\end{tabular}
\smallskip

\textbf{Equation $\leftrightarrow$ quantity map:}

\smallskip
\begin{tabular}{@{}p{2.8cm}p{3.8cm}p{6.2cm}@{}}
\toprule
Paper eq.                                    & Symbol                          & Quantity \\
\midrule
\eqref{eq:zhou-A-L}                          & $A$, $\overline L$              & risk-free annuity, DF-weighted Libor \\
\eqref{eq:zhou-asw}                          & $\Sasw$                         & par asset swap spread (two forms) \\
\eqref{eq:zhou-W}                            & $W$                             & default-neutral CDS notional ratio \\
\eqref{eq:zhou-pv-loss}                      & $\PVone^{\text{Risky}}$, $\Loss$ & risky annuity, default-leg integral \\
\eqref{eq:zhou-cds-standard}                 & $\Scds$ (standard)              & textbook CDS fair-spread identity \\
\eqref{eq:zhou-bond-decomp}                  & bond consistency                & bond price as coupon + recovery + survival-discounted notional \\
\eqref{eq:zhou-cds-explicit}--\eqref{eq:zhou-cds-simplified} & $\Scds$ (bond-implied) & explicit and simplified forms \\
\eqref{eq:zhou-basis}                        & $\Basis$                        & 3-term decomposition of CDS-bond basis \\
\bottomrule
\end{tabular}
\smallskip

\textbf{Invariants and identities to verify:}
\begin{enumerate}[noitemsep]
  \item \emph{Two ASW forms agree}: $(M - (1-D))/A = C - \overline L - D/A$
        identically (given the definitions of $M, A, \overline L$);
        the second form is the first rewritten using
        $M = \sum C\, \Delta T_{k-1}\, DF(T_k) + DF(T_N) = C \cdot A + 1$
        for a par bond on the risk-free curve.
  \item \emph{Default-neutral notional}: with CDS notional
        $(1 - D)/(1 - R)$ and recovery $R$, the default payment is
        $(1 - D)$, exactly matching the loan repayment. So the
        combined position has zero default exposure by construction.
  \item \emph{Consistency theorem}: under \eqref{eq:zhou-bond-decomp},
        the explicit \eqref{eq:zhou-cds-explicit} equals the
        textbook \eqref{eq:zhou-cds-standard}.
  \item \emph{Reproduction of Table 1}: at the calibrated
        $\lambda$ given in row 5, $\Scds, \Sasw, \Basis$ in the upper
        rows should be reproduced exactly. The three basis-decomposition
        rows must sum (with signs as in \eqref{eq:zhou-basis}) to
        the $\Basis$ row above.
  \item \emph{Curve-slope $\to$ basis sign}: in Table 2 (flat curve),
        the first basis term $\overline L - \overline L^{\text{Risky}} = 0$
        identically; the basis becomes (mostly) negative for low $D$
        and turns positive only as the off-par term dominates.
  \item \emph{Par bond basis $\neq 0$}: at $D = 0$, $\Basis = +0.01\%$
        on the upward-sloping curve and $-0.06\%$ on the flat curve
        --- showing that even par bonds have non-zero basis due to
        the accrual-asymmetry term.
  \item \emph{Monotonicity}: $\Scds$, $\Sasw$, $\lambda$, $\Basis$
        are all increasing in $D$ (Table 1). $\Scds$ grows faster
        than $\Sasw$, which is why $\Basis$ widens for deep-discount
        bonds.
  \item \emph{$W$ monotonicity in $D$}: $W$ decreases in $D$ for fixed
        $R$; at $D = 0$, $W = 1$. Check: $W = (1 - D / (1-R)) / (1 - D)$,
        differentiated and shown to be negative for $R < 1$.
\end{enumerate}
\end{validation}

% =====================================================================
\section{Glossary of symbols (paper's notation)}
\begin{tabular}{@{}p{3cm}p{12cm}@{}}
\toprule
$T_0, T_k$        & Today; bond / ASW payment dates \\
$\Delta T_k$      & Year fraction of period $(T_{k-1}, T_k)$ \\
$C$               & Bond fixed coupon (paper examples: 7\% semiannual) \\
$M$               & Bond risk-free price \\
$1 - D$           & Bond market dirty price; $D$ is the discount \\
$R$               & Expected recovery rate \\
$W$               & CDS-notional-to-bond-price ratio for default neutrality \\
$L_{k-1}$         & Forward Libor for period $(T_{k-1}, T_k)$ at $T_0$ \\
$DF(T_k)$         & Risk-free discount factor at $T_0$ \\
$P(\tau > T_k)$   & Survival probability from $T_0$ to $T_k$ \\
$\tau$            & Default time \\
$\lambda$         & Hazard rate (constant in \S2.3.1) \\
$A$               & Risk-free annuity \\
$\PVone$          & Risky annuity (uses $DF$ and $P$) \\
$\PVone^{\text{Risky}}$ & Survival-weighted risky annuity (paper's notation) \\
$\overline L$     & $A$-weighted average forward Libor \\
$\overline L^{\,\text{Risky}}$ & Risky-annuity-weighted average forward Libor \\
$\Loss$           & Risky default-leg integral, trapezoidal-approximated \\
$\Sasw$           & Par asset swap spread \\
$\Scds$           & Bond-implied CDS spread \\
$\Basis$          & $\Scds - \Sasw$, CDS-bond basis \\
\bottomrule
\end{tabular}

% =====================================================================
\begin{todobox}
\begin{itemize}[noitemsep]
  \item Reproduce Table 1 and Table 2 on the 10Y, 7\%-coupon bond
        with $R = 40\%$, using the 16-July-2008 Bloomberg swap zero
        curve. Match across the seven $D$ levels.
  \item Verify the 3-term basis decomposition sums to the $\Basis$
        row in both tables. The third row (off-par term) is given
        as $D / \PVone^{\text{Risky}} - D / A$, which differs by a
        sign convention from \eqref{eq:zhou-basis}; check.
  \item Implement \S2.3.2 bootstrap: given a bond price term
        structure on a single issuer, recover a piecewise-constant
        $\lambda$ from \eqref{eq:zhou-bond-decomp}.
  \item Implement \S2.3.3: given both market CDS spreads and bond
        prices on the same issuer, jointly solve for $\lambda(T)$
        and $R(T)$. Compare implied $R(T)$ to any available recovery
        swap quotes.
  \item Modern (OIS) update: re-do Table 1 with OIS discounting in
        place of Libor; quantify the impact on $\Scds$, $\Sasw$, and
        $\Basis$.
  \item Test the par-bond basis identity at $D = 0$ across various
        curve shapes (flat, upward, inverted, humped) to verify the
        sign predictions of the three-term decomposition.
\end{itemize}
\end{todobox}

\end{document}
